% Part: second-order-logic
% Chapter: syntax-and-semantics
% Section: expressive-power

\documentclass[../../../include/open-logic-section]{subfiles}

\begin{document}

\olfileid{sol}{syn}{exp}
\olsection{د بيان ځواک}

\begin{explain}
پر دويمې درجې متغيرونو کميت ټاکنه د لومړۍ درجې منطق په پرتله د ژبې د
بيان ځواک ډېر زياتوي۔ د دويمې درجې وجودي کميت ټاکنه پرېږدي چې ووايو
د ټاکلو خاصيتونو لرونکې تابعې يا اړيکې شته۔ په لومړۍ درجې منطق کښې
د دې يوازينۍ لار د همدې موخې دپاره د يو غيرمنطقي سمبول، يعنې
!!a{function} يا !!{predicate}، ټاکل دي۔ د دويمې درجې عمومي کميت
ټاکنه پرېږدي چې ووايو د !!{domain} ټول فرعي سټونه، پر هغې ټولې
اړيکې، يا له !!{domain} څخه هغې ته ټولې تابعې يو خاصيت لري۔ په
لومړۍ درجې منطق کښې يوازې دا ويلے شو چې د ژبې يوه غيرمنطقي سمبول ته
ګومارل شوے فرعي سټ، اړيکه يا تابع هغه خاصيت لري۔ د دې خاصيت په
ټاکلو کښې هم د دويمې درجې کميت ټاکنه کارولے شو، که وايو چې ځينې
فرعي سټونه، اړيکې او تابعې يې لري، او که وايو چې ټول يې لري۔ دا
موږ ته اجازه راکوي چې داسې اړيکې تعريف کړو چې په لومړۍ درجې منطق
کښې نۀ تعريفېږي، او د ساحې داسې خاصيتونه بيان کړو چې هلته نۀ شي
بيانېدلے۔
\end{explain}

\begin{defn}
که $\Struct{M}$ د ژبې~$\Lang{L}$ يو !!a{structure} وي، اړيکه
$R \subseteq \Domain{M}^2$ هله په~$\Lang{L}$ کښې \emph{تعريفېدونکې}
ده چې داسې !!{formula}~$!A_R(x, y)$ وي چې يوازې $x$ او~$y$ يې ازاد
متغيرونه وي، او $R(a, b)$، يعنې $\tuple{a,b} \in R$، هله او يوازې
هله برقرار وي چې د $s(x) = a$ او $s(y) = b$ دپاره
$\Sat{M}{!A_R(x, y)}[s]$ وي۔
\end{defn}

\begin{ex}
په لومړۍ درجې منطق کښې د عينيت اړيکه~$\Id{\Domain{M}}$، يعنې
$\Setabs{\tuple{a,a}}{a \in
  \Domain{M}}$، د $\eq[x][y]$ په فارمولا تعريفولے شو۔ په دويمې
درجې منطق کښې دا اړيکه \emph{له~$\eq$ پرته} هم تعريفولے شو۔ ځکه
چې که $a$ او $b$ د~$\Domain{M}$ يو شان !!{element} وي، د
~$\Domain{M}$ د هماغو فرعي سټونو !!{element}s دي (ځکه سټونه د خپلو
!!{element}s له مخې ټاکل کېږي)۔ برعکس، که $a$ او $b$ بېل وي، په
هماغو فرعي سټونو کښې !!{element}s نۀ دي: مثلاً که $a \neq b$ وي، $a \in
\{a\}$ خو $b \notin \{a\}$۔ نو ``د~$\Domain{M}$ په هماغو
فرعي سټونو کښې د !!{element}s کېدل'' داسې اړيکه ده چې پر $a$ او
$b$ هله او يوازې هله برقرار وي چې $a = b$۔ دا اړيکه په دويمې درجې
منطق کښې بيانېدے شي، ځکه د~$\Domain{M}$ پر ټولو فرعي سټونو کميت
ټاکلے شو۔ نو لاندې !!{formula} $\Id{\Domain{M}}$ تعريفوي:
\[
\lforall[X][(X(x) \liff X(y))]
\]
\end{ex}

\begin{prob}
وښيئ چې $\lforall[X][(X(x) \lif X(y))]$ (پام: $\lif$،
نۀ~$\liff$!) $\Id{\Domain{M}}$ تعريفوي۔
\end{prob}

\begin{ex}
که $R$ دوه ځايه !!{predicate} وي، $\Assign{R}{M}$ پر~$\Domain{M}$
دوه ځايه اړيکه ده۔ ښايي لږ مغشوشوونکې وي، خو $R$ به هم د
!!{predicate}~$R$ دپاره او هم د اړيکې~$\Assign{R}{M}$ پخپله دپاره
کاروو۔ د~$R$
\emph{متعدي تړون}~$R^*$ هغه اړيکه ده چې د $a$ او $b$ تر منځ هله
برقراره وي چې د ځينو $c_1$، \dots، $c_k$ دپاره $R(a,c_1)$،
$R(c_1, c_2)$، \dots، $R(c_k,b)$ برقرار وي۔ په دې کښې د $k = 0$
حالت هم شامل دے: که $R(a,b)$ وي، $R^*(a,b)$ هم وي۔ نو
$R \subseteq R^*$۔ په حقيقت کښې $R^*$ تر ټولو کوچنۍ اړيکه ده چې
~$R$ رانغاړي او متعدي ده۔ په دويمې درجې منطق کښې ويلے شو چې $X$
هغه متعدي اړيکه ده چې~$R$ رانغاړي:
\begin{multline*}
  !B_R(X) \ident \lforall[x][\lforall[y][(R(x,y) \lif X(x, y))]] \land {}\\
\lforall[x][\lforall[y][\lforall[z][((X(x,y) \land X(y,z)) \lif X(x,
      z))]]].
\end{multline*}
لومړے عطف وايي چې $R \subseteq X$، او دويم وايي چې $X$ متعدي ده۔

دا ويل چې $X$ د دې ډول تر ټولو کوچنۍ اړيکه ده، معنا يې دا ده چې
دا پخپله په هرې داسې اړيکې کښې شامله ده چې $R$ رانغاړي او متعدي ده۔
نو د~$R$ متعدي تړون په لاندې !!{formula} تعريفولے شو:
\[
R^*(X) \ident !B_R(X) \land \lforall[Y][(!B_R(Y) \lif
  \lforall[x][\lforall[y][(X(x, y) \lif Y(x,y))]])].
\]
$\Sat{M}{R^*(X)}[s]$ هله او يوازې هله وي چې $s(X) = R^*$۔ د~$R$
متعدي تړون په لومړۍ درجې منطق کښې نۀ شي بيانېدلے۔
\end{ex}

\end{document}
